\section{Efficient Inference: Model Compression}

% --- SLIDE 1: WHEN THE MODEL IS THE BOTTLENECK ---
\begin{frame}{When the Model Is the Bottleneck}

    \begin{itemize}
        \item The latency budget showed that for a typical web service the model is the \emph{smallest}
              term. That flips when the model runs \textbf{on the device} (a camera, a phone, a vehicle),
              when it is a \textbf{large network} served at high volume, or when \textbf{memory or energy}
              is the constraint.
        \item Then you make the model itself cheaper. Three levers, often combined:
    \end{itemize}

    \vspace{0.2em}
    \centering
    \renewcommand{\arraystretch}{1.3}
    {\scriptsize
    \begin{tabular}{p{2.0cm}|p{3.9cm}|p{3.3cm}|p{2.5cm}}
        \toprule
                              & \textbf{What changes} & \textbf{What you gain} & \textbf{Retraining?} \\
        \midrule
        \textbf{Quantisation} & Fewer bits per weight and activation (FP32 $\rightarrow$ INT8)
                              & $4\times$ smaller weights; faster where INT8 kernels exist
                              & None, or a fine-tune (QAT) \\
        \textbf{Pruning}      & Removes weights, channels or layers
                              & Smaller; faster only if the pruning is \emph{structured}
                              & Usually a fine-tune \\
        \textbf{Distillation} & Trains a small \emph{student} to mimic a large \emph{teacher}
                              & Any smaller architecture you choose
                              & A full training run \\
        \bottomrule
    \end{tabular}}

    \vspace{0.5em}
    \begin{block}{Rule}
        Measure on the \textbf{target hardware}, before and after. A compression that does not move p95
        latency, memory or energy on the device is cost without benefit.
    \end{block}

\end{frame}
